% 10-pairing.tex -- the half-complex pairing, map against visibility.
% Input by mostlyharmless.tex.  No preamble here.

\section{The coefficient pairing}
\label{sec:pairing}

A map and its visibility are expanded with the \emph{same} stored coefficients, and by the
same pairing.  This section collects that correspondence in one place, because it is used by
every later section, it is easy to get wrong, and the two bases do not treat negative orders
the same way.  Nothing here is new: it assembles results of
sections~\ref{sec:harmonics}, \ref{sec:wigner} and \ref{sec:visibility}.

\subsection{The map}
\label{sec:pairing-map}

The stored coefficients are those of order $m \ge 0$, and the map is real, so the negative
orders are determined rather than stored:
\begin{equation}
  a_{l,-m} \;=\; (-1)^m\,\overline{a_{lm}} .
  \label{eq:pairing-reality}
\end{equation}
Expanding over all orders and folding the negative ones back onto the stored half gives
\begin{equation}
  I(n) \;=\; \sum_{l}\Bigl[\, a_{l0}\,Y_l^0(n)
    \;+\; \sum_{m=1}^{l}\Bigl( a_{lm}\,Y_l^m(n)
    \;+\; (-1)^m\,\overline{a_{lm}}\,Y_l^{-m}(n) \Bigr) \Bigr] ,
  \label{eq:pairing-map}
\end{equation}
and because the harmonics obey
\begin{equation}
  Y_l^{-m} \;=\; (-1)^m\,\overline{Y_l^m} ,
  \label{eq:pairing-y-negative}
\end{equation}
the two terms are conjugates of one another and
\begin{equation}
  I(n) \;=\; \sum_{l}\Bigl[\, a_{l0}\,Y_l^0(n)
    \;+\; \sum_{m=1}^{l} 2\,\mathrm{Re}\bigl( a_{lm}\,Y_l^m(n) \bigr) \Bigr] .
  \label{eq:pairing-map-real}
\end{equation}

\subsection{The visibility}
\label{sec:pairing-visibility}

The same coefficients, the same pairing, and the same stored half give
\begin{equation}
  V(u,v) \;=\; \sum_{l}\Bigl[\, a_{l0}\,V_l^0(u,v)
    \;+\; \sum_{m=1}^{l}\Bigl( a_{lm}\,V_l^m(u,v)
    \;+\; (-1)^m\,\overline{a_{lm}}\,V_l^{-m}(u,v) \Bigr) \Bigr] .
  \label{eq:pairing-visibility}
\end{equation}
Here, however, the negative orders obey a different rule,
\begin{equation}
  V_l^{-m} \;=\; \overline{V_l^m} ,
  \label{eq:pairing-v-negative}
\end{equation}
so the second term is \emph{not} the conjugate of the first, and the $(-1)^m$ survives.  It
combines with the mode's own phase to leave
\begin{equation}
  a_{lm}V_l^m + (-1)^m\overline{a_{lm}}\,\overline{V_l^m}
  \;=\; 4\pi(-i)^m N_{lm}R_{lm}(q)\,\mathrm{Re}\bigl(a_{lm}e^{im\psi}\bigr) ,
  \label{eq:pairing-contribution}
\end{equation}
which is the contribution quoted in section~\ref{sec:visibility}.  So
\begin{equation}
  V(u,v) \;=\; \sum_{l}\Bigl[\, a_{l0}\,V_l^0(u,v)
    \;+\; \sum_{m=1}^{l} 4\pi(-i)^m N_{lm}R_{lm}(q)\,
      \mathrm{Re}\bigl( a_{lm}\,e^{im\psi} \bigr) \Bigr] .
  \label{eq:pairing-visibility-real}
\end{equation}

\subsection{Why the two negative orders differ}
\label{sec:pairing-why}

Equation~\eqref{eq:pairing-y-negative} is a definition of the basis.  For the visibility the
instruction ``take the mode at order $-m$'' contains a second operation: conjugating a
harmonic also conjugates its phase, and a plane wave whose phase is conjugated is the wave at
the opposite baseline.  What saves the derivation is
\begin{equation}
  V_l^m(-u,-v) \;=\; (-1)^m\,V_l^m(u,v) ,
  \label{eq:pairing-baseline-sign}
\end{equation}
which follows from $\psi \to \psi + \pi$ in \eqref{eq:mode-visibility}: the order-$m$ mode
picks up $e^{im\pi} = (-1)^m$.  Combined with the conjugation, the two signs cancel and leave
plain conjugation, \eqref{eq:pairing-v-negative}.  That cancellation is the whole reason the
map carries a $(-1)^m$ and the visibility does not, and the reason the mode sum shows a
$(-i)^m$ phase at all.

\subsection{What the pairing buys}
\label{sec:pairing-buys}

Three consequences, each used constantly:

\begin{itemize}
  \item only $m \ge 0$ is stored, halving the coefficient set, and half of what remains is
    real ($a_{l0}$);
  \item each stored coefficient costs exactly one mode evaluation, and the mode is a closed
    form (section~\ref{sec:odd} or~\ref{sec:even}), so the cost of a visibility is set by
    the number of modes and not by any sampling;
  \item the pairing is unchanged by rotation, since \eqref{eq:coefficient-rotation} delivers
    observer-frame coefficients in the same half-complex form.  A rotated map therefore
    needs no re-derivation, only a new coefficient set.
\end{itemize}

Identity \eqref{eq:pairing-baseline-sign} is \textbf{certificate \vtag{23}}: it is checked
directly against \eqref{eq:mode-visibility} at the two baselines, with no coefficients and no
pairing involved, so the identity is isolated rather than inferred.  That isolation is the
point.  The Hermitian symmetry of a real map, $V(-u,-v) = \overline{V(u,v)}$, also depends on
it, but only in combination with the pairing, so a sign error common to both would satisfy
the symmetry and still be wrong.  Certificate \vtag{23} would fail instead, and it would fail
loudly if the sense of the azimuth were ever changed again --- which it has been.

\begin{center}
\begin{tabular}{l l l}
\hline
tag & what is checked & kind \\
\hline
\vtag{23} & the mode at the opposite baseline, against \eqref{eq:mode-visibility} & algebraic \\
\hline
\end{tabular}
\end{center}
